\paragraph{F.44 (Open). Delimitation beyond the proved reconstruction ranges.}
The requested base-field delimitation now includes the following
proved positive ranges: infinite finitely generated fields by F.5,
sub-$p$-adic fields and the specified general bases by F.38, and
every algebraic subfield of $F\mathbb Q^{\mathrm{ab}}$, for a
number field $F$, by F.47, and, more generally, every subfield of
a finitely generated extension of $\mathbb Q^{\mathrm{ab}}$ by
F.50. These assertions include arbitrary cyclotomic extensions of
number fields and fields finitely generated over them. In particular, the test
over $\mathbb Q(\mu_{2^\infty})$ considered during this work is
settled, including full site realization. F.46 records its initial
point-function proof and F.47 supplies the additional argument
needed for the full geometric morphism. The proposed unrestricted
fullness is false by F.5. F.49 strengthens the counterexamples
F.6 and F.37: pinned fullness fails over every algebraic extension
of a finite field, every algebraically closed field, and every
infinite perfect field with nontrivial centre of its absolute
Galois group, including every real closed field.
Faithfulness of actual reduced-scheme maps remains
the all-characteristic theorem F.4.

The exact general criterion is Carlson--Haine--Wolf,
Theorem~4.18 \cite{CHW2024}: a field has pinned reconstruction if
and only if every pinned morphism from a connected regular affine
scheme with a rational point to $(\mathbb G_m)_{\mathrm{et}}$
has the geometric-point values of a regular unit (on the absolute
weak normalization in positive characteristic). Its implication to
full site realization is proved through all sliced etale objects,
as explained in F.47. Their Theorem~6.24 further supplies the
positive range of amenable fields admitting a nontrivial discrete
valuation. Amenability means that rational points of every torus
have no nonidentity prime-to-characteristic divisible element and
that the prime-to-characteristic Tate module of the Picard group of
every regular finite-type scheme vanishes (Definition~5.15).
This is an exact verified sufficient condition, not a claim that
every field outside it is a failure or a literature-wide open case.

A concrete test not decided by this return is
\[
 k_{\mathrm{solv}}=\mathbb Q^{\mathrm{solv}},\qquad
 \operatorname{Hom}_{k_{\mathrm{solv}}}(\mathbb G_m,\mathbb G_m)
 \longrightarrow
 \pi_0\operatorname{Hom}^{\mathrm{pin}}_{B G_{k_{\mathrm{solv}}}}
       ((\mathbb G_m)_{\mathrm{et}},(\mathbb G_m)_{\mathrm{et}}),
 \tag{F44}
\]
where $\mathbb Q^{\mathrm{solv}}$ is the maximal solvable algebraic
extension in a fixed algebraic closure. The map is injective by
F.4; its surjectivity is not asserted. Finite-type dimension-zero
schemes, after universal-homeomorphism localization, are finite
etale and their relative Galois topoi already determine their maps.
Thus (F44) is a first positive-dimensional test of the remaining
criterion, rather than a claim of a smallest base field or of
literature-wide minimality. Constant morphisms and the general
nonconstant, possibly noninvertible, locus are part of the question;
the unresolved realization issue concerns the latter locus.

Here are the precise failed extensions of the proof, rather than
an appeal to the absence of a cited theorem. The field
$k_{\mathrm{solv}}$ is radically closed. Indeed, an element lies
in a finite solvable Galois extension $E/\mathbb Q$. Adjoin the
needed roots of unity and then $n$th roots of all its conjugates;
the latter extension is abelian over the enlarged $E$, so its
normal closure is still solvable over $\mathbb Q$. Consequently
\[
       D(k_{\mathrm{solv}}^*)=k_{\mathrm{solv}}^*.
\]
The height argument used in F.47 cannot reduce these elements to
roots of unity: this field contains $2^{1/n}$ for every $n$, of
positive heights tending to zero. The same observation shows that
it cannot be contained in $F\mathbb Q^{\mathrm{ab}}$ for any
fixed number field $F$, by the verified Amoroso--Zannier bound.
It is also outside the broader range F.50: the algebraic constants
in a finitely generated extension of $\mathbb Q^{\mathrm{ab}}$
form a finite extension of that field, to which the same height
bound applies.
It is not amenable. It is not sub-$p$-adic either, since it contains
all roots of unity whereas the algebraic constants in a finitely
generated extension of $\mathbb Q_p$ form a finite extension and
have only finitely many roots of unity.

The Picard input of F.47 also fails here. For an elliptic curve
$E/\mathbb Q$, every $2^r$-division field is solvable: its Galois
group is a subgroup of
$\operatorname{GL}_2(\mathbb Z/2^r)$, an extension of
$\operatorname{GL}_2(\mathbb F_2)\simeq S_3$ by a $2$-group.
Hence $E[2^\infty]\subset E(k_{\mathrm{solv}})$, providing
nonzero infinitely divisible points and a nonzero Tate module.
Neither the integral-residue argument followed by Picard
divisibility nor the translated-function proof can therefore be
copied with its hypotheses unchanged. Kummer classes at the
$k_{\mathrm{solv}}$-points already forget all nonzero constants.

The relative birational isomorphism results checked in
Tsujimura's Theorem~4.7 and Corollary~4.10 \cite{Tsujimura2022}
do not supply this noninvertible realization step; the requisite
AVKF hypotheses also fail in this example. Descent of a scheme to
a number field does not imply descent of an arbitrary geometric
morphism on the whole etale site. No exotic pinned morphism in
(F44) has been constructed here, and no argument using the
nonabelian finite covers has yet eliminated one. The next step
would have to exploit those nonsolvable-cover compatibilities,
rather than the Kummer and Picard invariants just shown to vanish
or to have the wrong divisibility. This is the exact residual
test and obstruction left by this return; no broader assertion
about the status of all other fields is made.

A further literature check does not settle (F44).
Mochizuki's \emph{Global Solvably Closed Anabelian Geometry}
\cite{Mochizuki2006SolvablyClosed}, Definitions~1.1 and~2.3 and
Theorem~2.5, concerns quotient groups
$Q_F=\operatorname{Gal}(\widetilde F/F)$ for a solvably closed
extension of a number field, together with the required
$\Sigma$-compatible quotients of hyperbolic-curve fundamental
groups. Its conclusion reconstructs compatible isomorphisms.
The absolute group $G_{\mathbb Q^{\mathrm{solv}}}$ in (F44)
is instead a kernel in $G_{\mathbb Q}$, and (F44) asks about
all pinned morphisms, including noninvertible ones. Neither
changing the quotient to this kernel nor replacing the Isom
conclusion by Hom is justified by that theorem.

\medskip
\noindent\emph{Finite extensions do not repair these particular inputs.}
One possible approach would pass to a proper finite extension
$L/k_{\mathrm{solv}}$, which is Hilbertian by the Weissauer theorem
used explicitly in the proof of \cite[Proposition~2.5]{F44Jarden2011}.
Descent of a reconstruction theorem would indeed be valid here.
More generally, in characteristic zero pinned reconstruction
descends along a finite extension $L/k$. To prove this, apply the
criterion of \cite[Theorem~4.18]{CHW2024} to a connected regular
affine $X/k$ with a rational point and a pinned morphism to
$(\mathbb G_m)_{\mathrm{et}}$. The curve-free criterion applies
after finite etale base change to $L$, where it gives a unit
$f_L$ with the prescribed geometric-point values. On
$X_{L\otimes_kL}$ the two pullbacks of $f_L$ have the same values
at every geometric point, by the descent data of the original
morphism. This scheme is reduced, so the units agree.
Faithfully flat descent gives $f\in\Gamma(X,\mathcal O_X)^*$
with the required values over $k$; the same criterion proves
reconstruction. The source remains connected after base change:
regularity in characteristic zero makes it normal and smooth,
and its rational point forces its finite algebraic constant field
to be $k$.

The obstacle is therefore the arithmetic input over $L$, not this
descent argument. For \emph{every algebraic} extension
$L/k_{\mathrm{solv}}$, including each proper finite one, one has
\[
 k_{\mathrm{solv}}^*\subset D(L^*),\qquad
 E[2^\infty]\subset E(L).
\]
Thus $L^*/\mu_\infty$ contains the nonzero divisible subgroup
$k_{\mathrm{solv}}^*/\mu_\infty$ and cannot be free abelian;
the same elliptic curve still violates the AVKF and Tate-module
conditions. Nor can such an $L$ admit a nontrivial discrete
valuation. A homomorphism $v:L^*\to\mathbb Z$ arising from a
valuation vanishes on the divisible subgroup $k_{\mathrm{solv}}^*$.
For any element algebraic over this trivially valued field, its
minimal polynomial has all nonzero coefficients of value zero.
A positive or negative value of the element would make the
minimum value among its nonzero polynomial terms occur exactly
once, which is impossible. Hence $v$ vanishes on $L^*$.
Hilbertianity of a finite extension therefore does not restore
the multiplicative, Picard, or discrete-valuation hypotheses used
in the reconstruction proofs verified here. This eliminates that
specific route without producing an exotic pinned morphism or
deciding (F44).
